\documentclass[9pt,aspectratio=169,professionalfonts]{beamer}
\usetheme{Madrid}
\usecolortheme{default}
\setbeamertemplate{navigation symbols}{}

% Packages
\usepackage{amsmath}
\usepackage{amssymb}
% Algorithm packages not needed for this presentation
% \usepackage{algorithm}
% \usepackage{algorithmic}
\usepackage{listings}
\usepackage{graphicx}
\usepackage{tikz}
\usepackage{booktabs}
\usepackage{array}
% Simplified package setup for basic TeX Live
% \usepackage{multirow}
% \usepackage{colortbl}
% \usepackage{xcolor}
% \usepackage{pifont}

% Define checkmark symbol using basic LaTeX
\providecommand{\checkmark}{\textbf{✓}}

% Code listing style
\lstset{
    language=Python,
    basicstyle=\ttfamily\footnotesize,
    keywordstyle=\color{blue},
    commentstyle=\color{green!50!black},
    stringstyle=\color{red},
    showstringspaces=false,
    breaklines=true,
    frame=single,
    backgroundcolor=\color{gray!10}
}

% Title information
\title[Projection aléatoire pour graphes hétérogènes]{Projection aléatoire : embeddings efficaces pour graphes hétérogènes}
\subtitle{Présentation technique de 30 minutes}
\author{Your Name}
\institute{Your Institution}
\date{\today}

\begin{document}

% Title slide
\begin{frame}
\titlepage
\end{frame}

% Table of contents
\begin{frame}{Plan de la présentation (30 minutes)}
\tableofcontents
\end{frame}

\section{Définition du problème et motivation}

\begin{frame}{Le défi : Embedding d'un réseau bibliographique}
\begin{block}{Notre jeu de données}
\begin{itemize}
    \item 28,702 Auteurs
    \item 20 Conférences  
    \item 8,920 Termes
    \item Types de relations multiples : Auteur-Auteur, Auteur-Conférence, Auteur-Terme
\end{itemize}
\end{block}

\begin{block}{Objectif}
Apprendre des embeddings de dimension 256 qui capturent les relations sémantiques afin de prédire les collaborations
\end{block}

\begin{alertblock}{Pourquoi est-ce difficile ?}
\begin{itemize}
    \item Graphe hétérogène (plusieurs types de nœuds)
    \item Connectivité clairsemée
    \item Échelle : gérer efficacement plus de 28 000 nœuds
    \item Complexité sémantique : chaque type de relation a un sens différent
\end{itemize}
\end{alertblock}
\end{frame}

\begin{frame}{Approche traditionnelle – Marches aléatoires}
\begin{block}{Algorithme de marche aléatoire :}
\begin{itemize}
    \item Pour chaque nœud :
    \begin{itemize}
        \item Générer des marches aléatoires : $A \rightarrow C \rightarrow A \rightarrow T \rightarrow A \rightarrow \ldots$
        \item Considérer les marches comme des « phrases »
        \item Appliquer Word2Vec (Skip-gram) pour apprendre les embeddings
    \end{itemize}
\end{itemize}
\end{block}

\begin{alertblock}{Problèmes}
\begin{itemize}
    \item \textbf{Coût computationnel :} millions de marches nécessaires
    \item \textbf{Mémoire :} stockage de toutes les séquences
    \item \textbf{Biais d'échantillonnage :} les marches ratent la structure globale
    \item \textbf{Sensibilité aux hyperparamètres :} longueur et nombre de marches, etc.
\end{itemize}
\end{alertblock}

\textbf{Exemple :} Pour 28k auteurs, 10 marches \(\times\) 100 pas = 28M pas !
\end{frame}

\begin{frame}{Notre solution – Projection aléatoire : intuition}
\begin{block}{Idée clé : éviter les marches}
Instead of random walks, use random projections:
\begin{enumerate}
    \item \textbf{Create sparse random matrix} $\mathbf{R}$ (cheap to compute)
    \item \textbf{Apply graph operations directly:} $\mathbf{M} \mathbf{R}$, $\mathbf{M}^2 \mathbf{R}$, $\mathbf{M}^3 \mathbf{R}$
    \item \textbf{Combine features intelligently} with learned weights
\end{enumerate}
\end{block}

\begin{block}{Pourquoi ça marche}
Johnson-Lindenstrauss Lemma guarantees that random projections preserve distances in high-dimensional spaces.
\end{block}

\begin{exampleblock}{Analogy}
Instead of exploring a city by taking random walks, we take aerial photos from different angles and combine them!
\end{exampleblock}
\end{frame}

\section{Fondements théoriques}

\begin{frame}{Lemme de Johnson–Lindenstrauss}
\begin{theorem}[Johnson-Lindenstrauss Lemma (1984)]
For any set of $n$ points in high-dimensional space, there exists a mapping to $O(\log n)$ dimensions that preserves all pairwise distances within $(1\pm\varepsilon)$.
\end{theorem}

\begin{block}{Formulation formelle}
For points $\mathbf{x}_1, \mathbf{x}_2, \ldots, \mathbf{x}_n \in \mathbb{R}^d$, there exists random matrix $\mathbf{R} \in \mathbb{R}^{d \times k}$ where $k = O(\log n/\varepsilon^2)$ such that:

$$\boxed{(1-\varepsilon)\lVert\mathbf{x}_i - \mathbf{x}_j\rVert^2 \leq \lVert\mathbf{R}^T \mathbf{x}_i - \mathbf{R}^T \mathbf{x}_j\rVert^2 \leq (1+\varepsilon)\lVert\mathbf{x}_i - \mathbf{x}_j\rVert^2}$$
\end{block}

\begin{block}{Conséquences pour nous}
\begin{itemize}
    \item Random projections preserve neighborhood structure
    \item We can work in lower dimensions without losing information
    \item Provides theoretical justification for our approach
\end{itemize}
\end{block}
\end{frame}

\section{Algorithme et implémentation}

\begin{frame}{Étape 1 : Construction de la matrice de projection aléatoire}
\begin{block}{Dimensions matricielles}
$\mathbf{R} \in \mathbb{R}^{n \times d}$ (28,702 $\times$ 256 in our case)
\end{block}

\begin{block}{Motif de sparsité}
Each column has exactly $s=3$ non-zero entries:
$$\mathbf{R}_{ij} = \begin{cases}
\frac{+1}{\sqrt{s}} & \text{with probability } \frac{s}{n} \\
\frac{-1}{\sqrt{s}} & \text{with probability } \frac{s}{n} \\
0 & \text{otherwise}
\end{cases}$$
\end{block}

\begin{exampleblock}{Example for 6 authors, 4 dimensions, $s=2$}
$$\mathbf{R} = \begin{bmatrix}
0 & +0.7 & 0 & -0.7 \\
-0.7 & 0 & +0.7 & 0 \\
0 & 0 & -0.7 & +0.7 \\
+0.7 & -0.7 & 0 & 0 \\
0 & +0.7 & -0.7 & 0 \\
-0.7 & 0 & 0 & +0.7
\end{bmatrix}$$
\end{exampleblock}
\end{frame}

\begin{frame}{Étape 2 : Pondération par degré (paramètre \(\alpha\))}
\begin{alertblock}{Problem}
High-degree nodes (prolific authors) dominate embeddings
\end{alertblock}

\begin{block}{Solution}
Degree weighting with parameter $\alpha = -0.5$:
$$\boxed{\mathbf{R}' = \mathbf{D}^{\alpha} \mathbf{R}}$$
\end{block}

\begin{exampleblock}{Example with $\alpha = -0.5$}
\begin{align*}
\text{Author degrees:} &\quad [100, 10, 5, 50, 20, 8] \\
\mathbf{D}^{(-0.5)} &= \text{diag}([0.10, 0.32, 0.45, 0.14, 0.22, 0.35]) \\
\\
\text{Before weighting:} \\
A_1 \text{ (degree=100):} &\quad [0, +0.7, 0, -0.7] \\
A_2 \text{ (degree=10):} &\quad [-0.7, 0, +0.7, 0] \\
\\
\text{After weighting:} \\
A_1: &\quad [0, +0.07, 0, -0.07] \quad \text{\color{red} (Down-weighted)} \\
A_2: &\quad [-0.22, 0, +0.22, 0] \quad \text{\color{blue} (Up-weighted)}
\end{align*}
\end{exampleblock}
\end{frame}

\begin{frame}{Étape 3 : Matrices d'adjacence des méta-chemins}
\begin{block}{Les méta-chemins capturent les relations sémantiques}
\begin{align}
\text{\textbf{ACA} (Author-Conference-Author):} \quad &\mathbf{M}_{ACA} = \mathbf{A}_{AC} \mathbf{A}_{CA} \\
\text{\textbf{ATA} (Author-Term-Author):} \quad &\mathbf{M}_{ATA} = \mathbf{A}_{AT} \mathbf{A}_{TA} \\
\text{\textbf{AAA} (Author-Author-Author):} \quad &\mathbf{M}_{AAA} = \mathbf{A}_{AA} \mathbf{A}_{AA}
\end{align}
\end{block}

\begin{block}{Exemple de dimensions}
\begin{itemize}
    \item $\mathbf{A}_{AC}$: 28,702 $\times$ 20 (authors $\times$ conferences)
    \item $\mathbf{A}_{CA}$: 20 $\times$ 28,702 (conferences $\times$ authors)  
    \item $\mathbf{M}_{ACA}$: 28,702 $\times$ 28,702 (authors $\times$ authors)
\end{itemize}
\end{block}

\begin{block}{Interprétation}
$\mathbf{M}_{ACA}[i,j]$ = number of conferences authors $i$ and $j$ both published in
\end{block}
\end{frame}

\begin{frame}{Étape 4 : Normalisation (paramètre \(\beta\))}
\begin{alertblock}{Problem}
Some authors have many connections, others few
\end{alertblock}

\begin{block}{Solution}
Row normalization with $\beta = -1.0$:
$$\boxed{\tilde{\mathbf{M}} = \mathbf{D}^{\beta} \mathbf{M}}$$
\end{block}

\begin{exampleblock}{Example}
\begin{align*}
\text{Raw } \mathbf{M}_{ACA} \text{ row sums:} &\quad [50, 5, 100, 20] \quad \text{(conferences per author)} \\
\mathbf{D}^{(-1)} &= \text{diag}([1/50, 1/5, 1/100, 1/20]) \\
\\
\text{Before normalization:} \\
\text{Author}_1: &\quad [10, 5, 30, 5] \quad \text{(shared conferences)} \\
\text{Author}_2: &\quad [2, 0, 1, 2] \\
\\
\text{After normalization:} \\
\text{Author}_1: &\quad [0.2, 0.1, 0.6, 0.1] \quad \text{(probabilities)} \\
\text{Author}_2: &\quad [0.4, 0.0, 0.2, 0.4]
\end{align*}
\end{exampleblock}
\end{frame}

\begin{frame}{Étape 5 : Génération des caractéristiques – cœur de l'algorithme}
\begin{block}{Innovation clé}
Au lieu de calculer explicitement $\mathbf{M}^2$, $\mathbf{M}^3$, on itère :
\end{block}

\textbf{Algorithme des puissances de matrice (itératif) :}
\begin{itemize}
    \item $\mathbf{U}_0 = \mathbf{R}'$ \quad (Initial : matrice aléatoire pondérée)
    \item $\mathbf{U}_1 = \text{normalize}(\tilde{\mathbf{M}} \mathbf{U}_0)$ \quad (caractéristiques 1 saut)
    \item $\mathbf{U}_2 = \text{normalize}(\tilde{\mathbf{M}} \mathbf{U}_1)$ \quad (2-hop features)  
    \item $\mathbf{U}_3 = \text{normalize}(\tilde{\mathbf{M}} \mathbf{U}_2)$ \quad (3-hop features)
\end{itemize}

\begin{exampleblock}{Concrete Example (ACA meta-path)}
\begin{align*}
\mathbf{U}_0 &= \mathbf{R}' \in \mathbb{R}^{28,702 \times 256} \quad \text{(Random starting features)} \\
\mathbf{U}_1 &= \tilde{\mathbf{M}}_{ACA} \mathbf{U}_0 \quad \text{(Authors connected via 1 conference)} \\
\mathbf{U}_2 &= \tilde{\mathbf{M}}_{ACA} \mathbf{U}_1 \quad \text{(Authors connected via 2 conferences)} \\
\mathbf{U}_3 &= \tilde{\mathbf{M}}_{ACA} \mathbf{U}_2 \quad \text{(Authors connected via 3 conferences)}
\end{align*}
\end{exampleblock}

\begin{alertblock}{Memory Efficiency}
Never store $\tilde{\mathbf{M}}^2$ or $\tilde{\mathbf{M}}^3$ (would be 28K $\times$ 28K $\times$ 8 bytes = 6.4GB each!)
\end{alertblock}
\end{frame}

\begin{frame}{Étape 6 : Collecte multi-échelle des caractéristiques}
\begin{block}{Pour chaque couple méta-chemin $\times$ puissance}
\begin{align*}
\text{Meta-paths:} &\quad [AAA, ACA, ATA] \\
\text{Powers:} &\quad [1, 2, 3] \\
\text{Total features:} &\quad 3 \times 3 = 9 \text{ feature matrices} \\
\\
\text{Feature tensor:} &\quad \mathbf{F} \in \mathbb{R}^{9 \times 28,702 \times 256}
\end{align*}
\end{block}

\begin{block}{Interprétation des caractéristiques}
\begin{itemize}
    \item $\mathbf{F}[0]$: AAA$^1$ - Direct collaborations
    \item $\mathbf{F}[1]$: AAA$^2$ - 2-hop collaborations  
    \item $\mathbf{F}[2]$: AAA$^3$ - 3-hop collaborations
    \item $\mathbf{F}[3]$: ACA$^1$ - Shared conferences
    \item $\mathbf{F}[4]$: ACA$^2$ - Conference communities
    \item $\mathbf{F}[5]$: ACA$^3$ - Broader research areas
    \item $\mathbf{F}[6]$: ATA$^1$ - Shared keywords
    \item $\mathbf{F}[7]$: ATA$^2$ - Topic neighborhoods  
    \item $\mathbf{F}[8]$: ATA$^3$ - Research domains
\end{itemize}
\end{block}
\end{frame}

\begin{frame}{Étape 7 : Combinaison de caractéristiques apprenable}
\begin{alertblock}{Problem}
Which features are most important? Let the model decide!
\end{alertblock}

\begin{block}{Poids apprenables avec normalisation softmax}
\begin{align}
\boldsymbol{\theta} &\in \mathbb{R}^9 \quad \text{(Learnable parameters)} \\
\mathbf{w} &= \text{softmax}(\boldsymbol{\theta}) \quad \text{where } \sum_{i=0}^{8} w_i = 1 \\
\\
\mathbf{Z} &= \sum_{i=0}^{8} w_i \times \mathbf{F}[i] \quad \text{(Final embedding)}
\end{align}
\end{block}

\begin{exampleblock}{Example learned weights}
\begin{center}
\begin{tabular}{lll}
\toprule
AAA$^1$: 0.05 & AAA$^2$: 0.12 & AAA$^3$: 0.08 \\
\textbf{ACA$^1$: 0.25} & \textbf{ACA$^2$: 0.30} & \textbf{ACA$^3$: 0.15} \\
ATA$^1$: 0.02 & ATA$^2$: 0.02 & ATA$^3$: 0.01 \\
\bottomrule
\end{tabular}
\end{center}
\textbf{Insight:} Model learns that conference relationships are most predictive!
\end{exampleblock}
\end{frame}

\begin{frame}{Étape 8 : Prédiction de lien par distance}
\begin{block}{Score basé sur la distance}
For authors $i$ and $j$, get embeddings $\mathbf{z}_i, \mathbf{z}_j \in \mathbb{R}^{256}$:

\begin{align}
\text{distance}^2 &= \lVert\mathbf{z}_i - \mathbf{z}_j\rVert^2 = \sum_{k=1}^{256} (z_{i,k} - z_{j,k})^2 \\
\\
\text{probability} &= \sigma(\gamma - \lambda \times \text{distance}^2)
\end{align}
\end{block}

\begin{block}{Intuition}
\begin{itemize}
    \item Small distance $\rightarrow$ High probability of collaboration
    \item Large distance $\rightarrow$ Low probability of collaboration
\end{itemize}
\end{block}

\begin{block}{Paramètres apprenables}
\begin{itemize}
    \item $\gamma$ (intercept): Overall collaboration likelihood
    \item $\lambda$ (slope): How much distance matters
\end{itemize}
\end{block}
\end{frame}

\section{Optimisation et apprentissage}

\begin{frame}{Entraînement du modèle}
\begin{block}{Fonction de perte : deux termes}
$$\boxed{\mathcal{L} = \mathcal{L}_{BCE} + \lambda_{entropy} \times \mathcal{L}_{entropy}}$$
\end{block}

\begin{block}{1. Entropie croisée binaire (BCE)}
$$\mathcal{L}_{BCE} = -[y \log(\hat{p}) + (1-y) \log(1-\hat{p})]$$
\begin{itemize}
    \item Standard link prediction loss
    \item Weighted for class imbalance (3:1 negative:positive ratio)
\end{itemize}
\end{block}

\begin{block}{2. Régularisation entropique}
$$\mathcal{L}_{entropy} = -\sum_{i=1}^{9} w_i \log(w_i)$$
\begin{itemize}
    \item Prevents model from using only one feature type
    \item Encourages diverse feature usage
\end{itemize}
\end{block}

\begin{block}{Optimisation}
\begin{itemize}
    \item Adam optimizer (lr=0.01)
    \item Early stopping on validation AUC
    \item ReduceLROnPlateau scheduling
\end{itemize}
\end{block}
\end{frame}

\section{Optimisation des paramètres}

\begin{frame}{Optimisation des poids, pente et intercept}
\begin{block}{Principe}
\begin{enumerate}
  \item Les matrices de caractéristiques \(\mathbf{F}_f\) sont pré-calculées pour chaque méta-chemin et chaque puissance.
  \item Les poids \(w_f\) (issus d'un softmax sur les paramètres \(\theta_f\)) combinent ces matrices : \(\mathbf{E}=\sum_f w_f\,\mathbf{F}_f\).
  \item Pour chaque paire d'auteurs \(i,j\) : score \(s_{ij}=\gamma-\lambda\,\|\mathbf{e}_i-\mathbf{e}_j\|_2^2\).
  \item Perte totale : \( \mathcal L = \text{BCE}(s_{ij},y_{ij}) + \lambda_{\text{entropy}}\,\bigl(-\sum_f w_f\log w_f\bigr)\).
  \item Rétro-propagation et mise à jour des paramètres \(\theta_f,\lambda,\gamma\) via Adam ; arrêt anticipé sur l'AUC de validation.
\end{enumerate}
\end{block}
\end{frame}

\section{Avantages vs marches aléatoires}

\begin{frame}{Comparaison computationnelle}
\begin{table}
\centering
\begin{tabular}{lcc}
\toprule
\textbf{Aspect} & \textbf{Random Walks} & \textbf{FastRP} \\
& \textbf{(Node2Vec)} & \textbf{(Our Method)} \\
\midrule
Time Complexity & $O(rwl \times n)$ & $O(E \times d \times p)$ \\
Memory & $O(rwl \times n)$ & $O(d \times p)$ \\
Preprocessing & Generate walks & Sparse matrix ops \\
Scalability & Poor (quadratic) & Excellent (linear) \\
\bottomrule
\end{tabular}
\end{table}

\begin{block}{Where}
\begin{itemize}
    \item $r$ = walks per node (typically 10-50)
    \item $w$ = walk length (typically 80-100)  
    \item $l$ = training epochs (typically 100)
    \item $E$ = number of edges, $d$ = embedding dimension, $p$ = number of powers
\end{itemize}
\end{block}

\begin{alertblock}{Real Numbers for our dataset}
\begin{itemize}
    \item Node2Vec: $\sim$28K $\times$ 10 $\times$ 80 $\times$ 100 = 2.24B operations
    \item FastRP: $\sim$500K edges $\times$ 256 $\times$ 3 = 384M operations
    \item \textbf{$\sim$6x speedup!}
\end{itemize}
\end{alertblock}
\end{frame}

\begin{frame}{Bénéfices qualitatifs}
\begin{block}{Pourquoi les embeddings par projection aléatoire sont meilleurs}
\begin{enumerate}
    \item \textbf{Global Structure Preservation:}
    \begin{itemize}
        \item Random walks are local (limited by walk length)
        \item Matrix powers capture global connectivity patterns
    \end{itemize}
    
    \item \textbf{Multi-Scale Relationships:}
    \begin{itemize}
        \item Powers 1,2,3 capture immediate, community, and domain-level relationships
        \item Random walks only see fixed-length neighborhoods
    \end{itemize}
    
    \item \textbf{Deterministic Features:}
    \begin{itemize}
        \item No sampling variance
        \item Reproducible results
        \item Better for analysis and debugging
    \end{itemize}
    
    \item \textbf{Theoretical Guarantees:}
    \begin{itemize}
        \item Johnson-Lindenstrauss lemma provides distance preservation bounds
        \item Random walks have no such guarantees
    \end{itemize}
    
    \item \textbf{Heterogeneous Graph Support:}
    \begin{itemize}
        \item Natural handling of different relationship types
        \item Random walks struggle with heterogeneous graphs
    \end{itemize}
\end{enumerate}
\end{block}
\end{frame}

\section{Résultats et performances}

\begin{frame}{Résultats de prédiction de liens}
\begin{block}{Link Prediction Results}
\begin{itemize}
    \item \textbf{Training AUC:} 0.94
    \item \textbf{Validation AUC:} 0.94
    \item \textbf{Training Time:} 15 minutes (vs. hours for random walks)
\end{itemize}
\end{block}

\begin{block}{Importance des caractéristiques apprises}
\begin{center}
\begin{tabular}{lr}
\toprule
\textbf{Relationship Type} & \textbf{Weight} \\
\midrule
Conference relationships (ACA) & 70\% \\
Collaboration patterns (AAA) & 25\% \\
Topic similarity (ATA) & 5\% \\
\bottomrule
\end{tabular}
\end{center}
\end{block}

\begin{block}{Impact du filtrage de nœuds}
\begin{itemize}
    \item Original: 28,702 authors
    \item Filtered (min-degree=5): 27,821 authors
    \item Removed 881 low-degree "noise" nodes
    \item Improved clustering quality
\end{itemize}
\end{block}

\textbf{Scalability:} Linear scaling with graph size
\end{frame}

\begin{frame}{Exemple concret}
\begin{exampleblock}{Example: Predicting if Alice and Bob will collaborate}
\textbf{Step 1:} Alice (ID=1000), Bob (ID=1500)

\textbf{Step 2:} Extract their embeddings

\texttt{z\_alice = Z[1000, :] \# 256-dimensional vector}

\texttt{z\_bob = Z[1500, :] \# 256-dimensional vector}

\textbf{Step 3:} Compute distance

\texttt{distance\_sq = norm(z\_alice - z\_bob)**2}

\texttt{\# Result: 2.34}

\textbf{Step 4:} Predict probability

\texttt{gamma = 0.8 \# learned intercept}

\texttt{slope = 0.5 \# learned slope}

\texttt{probability = sigmoid(gamma - slope * 2.34)}

\texttt{\# Result: 0.41}

\textbf{Result:} 41\% chance Alice and Bob will collaborate

\textbf{How did we get z\textunderscore alice?} From weighted combination of 9 different relationship features!
\end{exampleblock}
\end{frame}

\begin{frame}{Résumé comparatif}
\begin{block}{Projection aléatoire vs méthodes par marches aléatoires}
\textbf{Advantages of Projection aléatoire:}
\begin{itemize}
    \item[\checkmark] \textbf{Speed:} 6x faster than Node2Vec
    \item[\checkmark] \textbf{Memory:} Constant memory usage
    \item[\checkmark] \textbf{Quality:} Better global structure capture
    \item[\checkmark] \textbf{Deterministic:} No sampling variance
    \item[\checkmark] \textbf{Scalable:} Linear time complexity
    \item[\checkmark] \textbf{Interpretable:} Feature weights show what matters
    \item[\checkmark] \textbf{Heterogeneous:} Natural multi-relation support
\end{itemize}

\textbf{Potential Disadvantages:}
\begin{itemize}
    \item[\textbf{!}] \textbf{Complexity:} More parameters to tune ($\alpha$, $\beta$, powers)
    \item[\textbf{!}] \textbf{Theory gap:} Less established than random walks
    \item[\textbf{!}] \textbf{Implementation:} Requires careful sparse matrix handling
\end{itemize}
\end{block}

\begin{alertblock}{Bottom Line}
Projection aléatoire provides superior speed-quality tradeoff for heterogeneous graphs
\end{alertblock}
\end{frame}

\section{Points essentiels}

\begin{frame}{Points à retenir}
\begin{block}{1. Core Insight}
Random projections + matrix powers = efficient graph embeddings
\end{block}

\begin{block}{2. Mathematical Foundation}
Johnson-Lindenstrauss lemma justifies random projection approach
\end{block}

\begin{block}{3. Engineering Innovation}
Iterative computation avoids memory explosion: $\mathbf{M} \circ (\mathbf{M} \circ \mathbf{R})$ not $\mathbf{M}^2 \circ \mathbf{R}$
\end{block}

\begin{block}{4. Performance Gains}
6x speedup with comparable/better quality vs. random walks
\end{block}

\begin{block}{5. Practical Impact}
Real-world heterogeneous graphs benefit from multi-relation modeling
\end{block}

\begin{alertblock}{The Big Picture}
Projection aléatoire makes graph embedding practical for large, complex, heterogeneous networks
\end{alertblock}
\end{frame}

\begin{frame}{Merci !}
\begin{block}{Code et implémentation}
\begin{itemize}
    \item Available on GitHub: [Your Repository]
    \item Python implementation with PyTorch
    \item Supports CPU, CUDA, and Apple Silicon (MPS)
\end{itemize}
\end{block}

\begin{block}{Fichiers clés}
\begin{itemize}
    \item \texttt{src/model.py}: Projection aléatoire implementation
    \item \texttt{src/main.py}: Training pipeline
    \item \texttt{scripts/plot\textunderscore umap.py}: Visualization tools
\end{itemize}
\end{block}

\begin{block}{Essayez vous-même}
\texttt{python src/main.py -epochs 30 -meta-paths AAA ACA ATA}
\end{block}

\begin{center}
\Huge \textbf{Questions?}
\end{center}
\end{frame}

% Appendix
\appendix

\begin{frame}{Appendix: Matrix Dimension Reference}
\begin{block}{Données d'entrée}
\begin{itemize}
    \item Authors: 28,702
    \item Conferences: 20  
    \item Terms: 8,920
\end{itemize}
\end{block}

\begin{block}{Matrices de relations}
\begin{itemize}
    \item $\mathbf{A}_{AA}$: 28,702 $\times$ 28,702 (sparse)
    \item $\mathbf{A}_{AC}$: 28,702 $\times$ 20
    \item $\mathbf{A}_{AT}$: 28,702 $\times$ 8,920
    \item $\mathbf{A}_{CA}$: 20 $\times$ 28,702
    \item $\mathbf{A}_{TA}$: 8,920 $\times$ 28,702
\end{itemize}
\end{block}

\begin{block}{Matrices des méta-chemins}
\begin{itemize}
    \item $\mathbf{M}_{AAA}$: 28,702 $\times$ 28,702
    \item $\mathbf{M}_{ACA}$: 28,702 $\times$ 28,702  
    \item $\mathbf{M}_{ATA}$: 28,702 $\times$ 28,702
\end{itemize}
\end{block}

\begin{block}{Matrices de caractéristiques}
\begin{itemize}
    \item Each $\mathbf{U}_i$: 28,702 $\times$ 256
    \item Feature tensor $\mathbf{F}$: 9 $\times$ 28,702 $\times$ 256
    \item Final embeddings $\mathbf{Z}$: 28,702 $\times$ 256
\end{itemize}
\end{block}
\end{frame}

\begin{frame}{Appendix: Memory Usage}
\begin{block}{Répartition mémoire}
\begin{itemize}
    \item Sparse matrices: $\sim$50MB
    \item Feature tensor: $\sim$2.5GB
    \item Total: $\sim$3GB (vs. 50GB+ for naive approach)
\end{itemize}
\end{block}

\begin{block}{Détails clés de l'implémentation}
\begin{itemize}
    \item Used SciPy sparse matrices throughout
    \item Never materialize full $\mathbf{M}^2$ matrices (would exceed memory)
    \item Iterative computation: $\mathbf{M} \circ (\mathbf{M} \circ \mathbf{R})$ instead of $\mathbf{M}^2 \circ \mathbf{R}$
    \item Disk cache for meta-path matrices
    \item GPU optimization with hybrid CPU/GPU approach
\end{itemize}
\end{block}
\end{frame}

\begin{frame}{Visualisation des embeddings}
\centering
\includegraphics[width=0.9\linewidth]{figures/embedding.png}
\end{frame}

\end{document} 